\documentclass[12pt]{article}

\usepackage{amsmath}
\usepackage{amssymb}
\usepackage{amsthm}
\usepackage{url}

\setlength{\textwidth}{6.5in}
\setlength{\oddsidemargin}{0in}
\setlength{\evensidemargin}{0in}
\setlength{\topmargin}{-0.25in}
\setlength{\textheight}{8.6in}

\theoremstyle{plain}
\newtheorem{theorem}{Theorem}
\newtheorem{lemma}[theorem]{Lemma}
\newtheorem{proposition}[theorem]{Proposition}
\newtheorem{corollary}[theorem]{Corollary}

\theoremstyle{definition}
\newtheorem{definition}[theorem]{Definition}

\theoremstyle{remark}
\newtheorem{remark}[theorem]{Remark}

\newcommand{\N}{\mathbb{N}}
\newcommand{\V}{\mathcal{V}}
\newcommand{\Sd}{\mathcal{S}}
\newcommand{\lo}[1]{\underline{#1}}
\newcommand{\hi}[1]{\overline{#1}}

\begin{document}

\begin{center}
{\LARGE\bf The greedy sequence avoiding 3-term arithmetic
progressions beyond its seed:\\[2mm] a proof of the conjectured
closed form of OEIS A092482}\\[8mm]
{\large Neal Patel}\\[1mm]
\texttt{neal@patel.codes}
\end{center}

\vskip 0.2in

\begin{abstract}
Sequence A092482 of the On-Line Encyclopedia of Integer Sequences
is the increasing sequence obtained greedily from $a(1)=1$,
$a(2)=2$, $a(3)=3$ by always appending the least integer that
creates no $3$-term arithmetic progression other than the initial
triple $(1,2,3)$.  The entry conjectures a closed form, recorded
there as ``conjectured and checked up to $n=512$.''  We prove it:
for $n\geq 1$,
$a(n+2)=1+2^{\lfloor\log_2 n\rfloor}+s(n)$, where $s(n)$ is the
value of the binary expansion of $n$ read in base $3$ (sequence
A005836).  The term set decomposes into blocks, each a translate
of an initial segment of the Stanley sequence A003278, and the
proof is the classical base-$2$/base-$3$ greedy argument applied
blockwise; the single new ingredient is a numerical coincidence
that makes each greedy jump land exactly at the start of the next
block.  Despite its seed, the sequence is not a Stanley sequence
in the sense of Odlyzko--Stanley, and it lies outside the
classification framework of Rolnick and Moy--Rolnick.  The proof
has been formalized and machine-checked in Lean~4 over Mathlib.
\end{abstract}

\section{Introduction}\label{sec:intro}

Entry A092482 of the On-Line Encyclopedia of Integer Sequences
(OEIS) \cite{OEIS} is titled ``Sequence contains no 3-term
arithmetic progression, other than its initial terms 1, 2, 3,''
and is defined in its comments by the greedy rule:

\begin{quote}
$a(1)=1$, $a(2)=2$, $a(3)=3$; $a(n)$ is least $k$ such that no
three terms of $a(1), a(2), \ldots, a(n-1), k$ form an arithmetic
progression, except for the first triple $(1,2,3)$.
\end{quote}

\noindent Its terms begin
\[
1,\ 2,\ 3,\ 6,\ 7,\ 14,\ 15,\ 17,\ 18,\ 36,\ 37,\ 39,\ 40,\ 45,\
46,\ 48,\ 49,\ 98,\ 99,\ \ldots
\]
Had the triple $(1,2,3)$ been forbidden rather than exempted, the
greedy rule starting from $1$ would produce the classical Stanley
sequence A003278 $=1,2,4,5,10,11,13,14,\ldots$ of Odlyzko and
Stanley \cite{OS1978}; the exemption changes the sequence from
the third term on, and, as we will see, changes its structure in
an essential way.

The formula field of the entry states, verbatim:

\begin{quote}
\texttt{For n > 2, a(n+2) = 1 + 2\textasciicircum
floor(log\_2(n)) + Sum\_\{k=1..n\} (3\textasciicircum
A007814(n) + 1)/2 = 1 + A053644(n) + A005836(n) (conjectured and
checked up to n=512).}
\end{quote}

\noindent The entry gives no proof and, as of this writing, still
labels the formula conjectural.  The unambiguous rendering of the
formula is the Mathematica program contributed to the entry by
Jean-Fran\c{c}ois Alcover on January 15, 2019, which compares the
first $512$ terms against
\[
b(n) \;=\; 1 + 2^{\,\ell-1} + \mathrm{FromDigits}(\mathrm{bits}(n-2),3),
\qquad n \geq 4,
\]
where $\mathrm{bits}(m)$ is the binary digit list of $m$ and
$\ell$ its length.  (The printed formula itself contains two
small typographical slips, confirmed computationally and
documented in Appendix~\ref{app:errata}; neither affects the
statement proved here.)

To state the result cleanly, for $n\geq 0$ write
$n=\sum_i b_i 2^i$ with $b_i\in\{0,1\}$ and define
\[
s(n) \;=\; \sum_i b_i 3^i,
\]
the binary expansion of $n$ read in base $3$.  Thus
$s(0),s(1),s(2),\ldots = 0,1,3,4,9,10,12,13,27,\ldots$ is the
sequence A005836 of numbers whose base-$3$ representation contains
no digit $2$, indexed from $0$.  Our main result is the conjectured
formula:

\begin{theorem}\label{thm:main}
For every $n\geq 1$,
\[
a(n+2) \;=\; 1 + 2^{\lfloor\log_2 n\rfloor} + s(n).
\]
\end{theorem}

\noindent The OEIS states the formula for $n>2$; it in fact holds
from $n=1$ on, and for $n=1,2$ it reproduces $a(3)=3$ and
$a(4)=6$.  Since $2^{\lfloor\log_2 n\rfloor}$ is A053644 and
$s(n)$ is A005836 (both read as $0$-indexed lists; see
Appendix~\ref{app:errata}), this is precisely the conjectured
closed form.  The summation form of the formula, with A007814
$=\nu_2$, the $2$-adic valuation, is Proposition~\ref{prop:sum}
below.

The engine of the proof is a block decomposition of the term set
(Theorem~\ref{thm:blocks}): apart from the two seed terms $1,2$,
the terms come in consecutive blocks, the $L$-th block being the
translate by $2^L+3^L+1$ of the first $2^L$ terms of A005836 ---
that is, of an initial segment of the (shifted) Stanley sequence
A003278.  Within a block, admissibility and minimality are the
classical base-$3$ digit arguments for A003278 going back to
\cite{OS1978}; what needs proof beyond that is the behavior at
block boundaries, which rests on the identity
\[
2\cdot\max(\text{block } L) \;=\; 2^{L+1}+3^{L+1}+1
\;=\; \min(\text{block } L+1):
\]
every integer strictly between two consecutive blocks is blocked
by an earlier arithmetic progression, and the first integer that
\emph{cannot} be blocked --- because doubling any earlier term
falls short of reaching past it --- is exactly the start of the
next block.  We are not aware of any prior proof of the closed
form; the OEIS entry, and the companion array A093682 which
collects conjectured closed forms for this family of greedy
sequences, both still describe the formulas as unproved.  We make
no claim of methodological novelty: the mechanism is the standard
base-$2$/base-$3$ greedy argument, applied blockwise.  In genre,
this note follows recent papers of Fried
\cite{Fried2024,Fried2026} proving conjectures stated in OEIS
entries.

Although its terms are built from translates of A003278, the
sequence A092482 itself is \emph{not} a Stanley sequence, because
a Stanley sequence requires a progression-free seed and
$\{1,2,3\}$ is a $3$-term progression; Section~\ref{sec:stanley}
makes this precise and explains why the classification machinery
for Stanley sequences does not apply.

The entire proof has been formalized in Lean~4 over the Mathlib
library and is machine-checked end to end;
Section~\ref{sec:lean} describes the formalization and what it
guarantees, including that the greedy sequence is defined in the
proof assistant directly from the OEIS rule, with no reference to
the closed form.

\section{The block decomposition}\label{sec:blocks}

Throughout, $\N=\{0,1,2,\ldots\}$.  Let
$\Sd=\{s(n):n\in\N\}=\{0,1,3,4,9,10,12,13,27,\ldots\}$ be the set
of naturals whose base-$3$ digits lie in $\{0,1\}$ (A005836).
The map $s$ satisfies $s(2n)=3s(n)$ and $s(2n+1)=3s(n)+1$, and is
strictly increasing.  We collect the digit facts we use; all are
elementary and classical.

\begin{lemma}\label{lem:digits}
Let $L\geq 0$.
\begin{enumerate}
\item[(i)] $s(2^L)=3^L$, and $s(2^L+r)=3^L+s(r)$ for $0\leq r<2^L$.
\item[(ii)] If $r<2^L$ then $2\,s(r)<3^L$.  Consequently
$\Sd\cap[0,3^L)=\{s(r):0\leq r<2^L\}$, a set with $2^L$ elements
and maximum $s(2^L-1)=(3^L-1)/2$.
\item[(iii)] (No progressions in $\Sd$.)  If $x,y,z\in\Sd$ and
$x+z=2y$, then $x=y=z$.
\end{enumerate}
\end{lemma}

\begin{proof}
(i) is immediate from the definition of $s$: prefixing the binary
digit $1$ prefixes the ternary digit $1$.  For (ii), the base-$3$
digits of $s(r)$ are supported on positions $<L$ and lie in
$\{0,1\}$, so $2\,s(r)$ has digits in $\{0,2\}$ on positions
$<L$, whence $2\,s(r)\leq 3^L-1$; the description of
$\Sd\cap[0,3^L)$ follows from (i) and strict monotonicity.  For
(iii): the base-$3$ addition $x+z$ is carry-free because each
digit sum is at most $2$, and each digit of $2y$ lies in
$\{0,2\}$.  Comparing digits, at every position the digits of $x$
and $z$ have an even sum, and since they lie in $\{0,1\}$ they
are equal.  Hence $x=z$, and then $y=x$.
\end{proof}

Part (iii) is the reason A003278 is the greedy progression-free
sequence: A003278$(n)=s(n-1)+1$, i.e., A003278 is the translate
$\Sd+1$ in increasing order \cite{OS1978}.

We also need the two digit witnesses that drive the classical
minimality argument.  For $u\in\N$ with base-$3$ digits $d_i$,
let
\[
\lo{u} = \text{the number with digits } (d_i \mapsto 0,1,0
\text{ for } d_i=0,1,2), \qquad
\hi{u} = \text{the number with digits } \min(d_i,1),
\]
i.e., $\lo{u}$ keeps the $1$-digits of $u$ and zeroes its
$2$-digits, while $\hi{u}$ caps every digit at $1$.

\begin{lemma}\label{lem:witness}
For every $u\in\N$: $\lo{u},\hi{u}\in\Sd$;
$\lo{u}\leq\hi{u}\leq u$; and
\[
u+\lo{u} \;=\; 2\,\hi{u}.
\]
If moreover $u\notin\Sd$, then $\hi{u}<u$ (and so also
$\lo{u}<u$).
\end{lemma}

\begin{proof}
The digit maps take values in $\{0,1\}$, so both witnesses lie in
$\Sd$, and both are digitwise dominated by $u$, hence $\leq u$.
The identity holds digit by digit with no carries:
$0+0=2\cdot 0$, $1+1=2\cdot 1$, $2+0=2\cdot 1$.  If $u\notin\Sd$
then some digit equals $2$ and is strictly lowered by capping, so
$\hi{u}<u$.
\end{proof}

\begin{corollary}[reflection]\label{cor:reflect}
For every $L\geq 0$ and every $u<3^L$ there exist
$t,t'\in\Sd\cap[0,3^L)$ with $u+t'=2t$.
\end{corollary}

\begin{proof}
If $u\in\Sd$ take $t=t'=u$; otherwise take $t=\hi{u}$,
$t'=\lo{u}$, both $\leq u<3^L$, and apply
Lemma~\ref{lem:witness}.
\end{proof}

Now the decomposition.  For $L\geq 0$ set
\[
C_L \;=\; 2^L+3^L+1, \qquad
T_L \;=\; \Sd\cap[0,3^L), \qquad
B_L \;=\; C_L+T_L \;=\; \{\,C_L+t : t\in T_L\,\},
\]
and define
\[
\V \;=\; \{1,2\}\;\cup\;\bigcup_{L\geq 0} B_L .
\]
The block $B_L$ has $2^L$ elements, starts at $C_L$
($C_0,C_1,C_2,\ldots=3,6,14,36,98,276,\ldots$), and by
Lemma~\ref{lem:digits}(ii) its maximum is
$C_L+(3^L-1)/2$.  Write
\[
P_L \;=\; \{1,2\}\cup B_0\cup\cdots\cup B_{L-1}
\]
for the seed together with the blocks of index $<L$.  The key
numerical coincidence is part (i) of:

\begin{lemma}\label{lem:double}
For every $L\geq 0$:
\begin{enumerate}
\item[(i)] $2\cdot\max B_L = C_{L+1}$;
\item[(ii)] the blocks are pairwise disjoint and ordered, with
$\max B_M<C_{M+1}$ for all $M$; consequently
$\V\cap[0,C_L)=P_L$;
\item[(iii)] if $L\geq 1$, every $p\in P_L$ satisfies
$2p\leq C_L$, with equality only for $p=\max B_{L-1}$.
\end{enumerate}
\end{lemma}

\begin{proof}
(i) is a computation:
\[
2\Bigl(C_L+\tfrac{3^L-1}{2}\Bigr)
= 2^{L+1}+2\cdot 3^L+2+3^L-1
= 2^{L+1}+3^{L+1}+1
= C_{L+1}.
\]
(ii) By (i), $\max B_M<2\max B_M=C_{M+1}$, and $C$ is strictly
increasing, so each block lies strictly below the start of the
next; the elements $1,2$ lie below $C_0=3$.  Hence
$\V\cap[0,C_L)=P_L$.  (iii) If $p\in B_M$ with $M\leq L-1$, then
$2p\leq 2\max B_M=C_{M+1}\leq C_L$, with equality forcing
$M=L-1$ and $p=\max B_{L-1}$; if $p\in\{1,2\}$, then
$2p\leq 4<6=C_1\leq C_L$.
\end{proof}

Our structural theorem, proved over the next three sections, is:

\begin{theorem}\label{thm:blocks}
The set of terms of A092482 is $\V$.
\end{theorem}

Theorem~\ref{thm:main} follows by counting: granting
Theorem~\ref{thm:blocks}, the terms in increasing order are $1$,
$2$, then the blocks in order, so for $L\geq 0$ and
$0\leq r<2^L$ the $(2^L+2+r)$-th term is $C_L+s(r)$.  Given
$n\geq 1$, write $n=2^L+r$ with $L=\lfloor\log_2 n\rfloor$ and
$0\leq r<2^L$; then by Lemma~\ref{lem:digits}(i)
\[
a(n+2) \;=\; C_L+s(r) \;=\; 1+2^L+\bigl(3^L+s(r)\bigr)
\;=\; 1+2^{\lfloor\log_2 n\rfloor}+s(n).
\]

\section{Admissibility: no progressions beyond the seed}
\label{sec:admissible}

\begin{theorem}\label{thm:admissible}
If $x<y<z$ are elements of $\V$ with $x+z=2y$, then
$(x,y,z)=(1,2,3)$.
\end{theorem}

\begin{proof}
Since $z\geq 3$, we have $z\in B_L$ for a (unique) $L$.

If $L=0$: $B_0=\{3\}$, so $z=3$ and $x<y<3$ force
$x,y\in P_0=\{1,2\}$, i.e., $y=2$, $x=1$, and indeed
$1+3=2\cdot 2$.

Now let $L\geq 1$; we show the progression is impossible.  Write
$z=C_L+t_z$ with $t_z\in T_L$.

\emph{First, $y\in B_L$.}  Suppose not.  Since $y<z<C_{L+1}$ and
$\V\cap[0,C_{L+1})=P_{L+1}=P_L\cup B_L$
(Lemma~\ref{lem:double}(ii)), we get $y\in P_L$, so
$2y\leq C_L$ by Lemma~\ref{lem:double}(iii); but then
$x+z=2y\leq C_L\leq z$ gives $x\leq 0$, a contradiction.  So
$y=C_L+t_y$ with $t_y\in T_L$.

\emph{Second, $x\in B_L$.}  Suppose not; as before $x\in P_L$
and $2x\leq C_L$.  From $x=2y-z=C_L+2t_y-t_z$ we get
$2x=2C_L+4t_y-2t_z$, so $2x\leq C_L$ forces
$C_L+4t_y\leq 2t_z$.  But $2t_z<3^L$ by
Lemma~\ref{lem:digits}(ii), while $C_L>3^L$ --- a contradiction.
So $x=C_L+t_x$ with $t_x\in T_L$.

Subtracting the translations, $t_x+t_z=2t_y$ with
$t_x,t_y,t_z\in\Sd$, so $t_x=t_z$ by
Lemma~\ref{lem:digits}(iii), contradicting $x<z$.
\end{proof}

\begin{remark}\label{rem:landing}
The equality case of Lemma~\ref{lem:double} explains why the
greedy sequence jumps from $\max B_L$ to exactly $C_{L+1}$.  To
\emph{block} an integer $m$ is to exhibit terms $x<y<m$ with
$x+m=2y$; any such $y$ satisfies $2y=x+m>m$.  Since every term
$y<C_{L+1}$ has $2y\leq C_{L+1}$
(Lemma~\ref{lem:double}(iii) at index $L+1$), the earlier terms
cannot block $C_{L+1}$ itself: even the extreme choice
$y=\max B_L$ would need $x=2y-C_{L+1}=0\notin\V$.  The next
section shows that, in contrast, everything strictly between
$\max B_L$ and $C_{L+1}$ \emph{is} blocked: $C_{L+1}$ is the
first integer past block $L$ that cannot be blocked, and the
greedy jump lands exactly there.
\end{remark}

\section{Minimality: every excluded integer is blocked}
\label{sec:minimal}

\begin{theorem}\label{thm:blocked}
Every integer $m>2$ with $m\notin\V$ admits $x,y\in\V$ with
$x<y<m$ and $x+m=2y$.
\end{theorem}

Within the span of a block this is the classical digit-witness
argument; across the gap between consecutive blocks it is the
following covering lemma, whose four cases tile the gap exactly.

\begin{lemma}[gap covering]\label{lem:gap}
For all $L\geq 0$ and all $\mu$ with $0\leq\mu<3^L+2^L$ there
exist $t\in T_L$ and $p\in P_L$ such that
\[
\mu+p \;=\; 2t+2^L+1 .
\]
\end{lemma}

\begin{proof}
Induction on $L$.

\emph{Base $L=0$.}  Here $\mu\in\{0,1\}$, $T_0=\{0\}$, and
$P_0=\{1,2\}$.  For $\mu=0$ take $t=0$, $p=2$; for $\mu=1$ take
$t=0$, $p=1$.  (These are the progressions $2,3,4$ and $1,3,5$
blocking the two integers $4,5$ between $B_0$ and $B_1$.)

\emph{Step $L\to L+1$.}  Let
$\mu<3^{L+1}+2^{L+1}=3\cdot 3^L+2\cdot 2^L$.  We use four cases,
whose ranges have lengths $2^L$, $3^L$, $3^L$, and $3^L+2^L$, and
therefore tile $[0,\,3\cdot 3^L+2\cdot 2^L)$ exactly.

\emph{Case A: $\mu<2^L$.}  Set $u=\mu+3^L-2^L$, so
$0\leq u<3^L$.  By Corollary~\ref{cor:reflect} there are
$t,t'\in T_L$ with $u+t'=2t$.  Take this $t\in T_L\subseteq
T_{L+1}$ and $p=C_L+t'\in B_L\subseteq P_{L+1}$.  Then
\[
\mu+p \;=\; (u+2^L-3^L)+(2^L+3^L+1+t')
\;=\; (u+t')+2^{L+1}+1 \;=\; 2t+2^{L+1}+1 .
\]

\emph{Case B: $2^L\leq\mu<2^L+3^L$.}  Set $\mu'=\mu-2^L<3^L$.
The induction hypothesis gives $t\in T_L$, $p\in P_L$ with
$\mu'+p=2t+2^L+1$; adding $2^L$ to both sides gives
$\mu+p=2t+2^{L+1}+1$, with $t\in T_{L+1}$ and $p\in P_{L+1}$.

\emph{Case C: $2^L+3^L\leq\mu<2^L+2\cdot 3^L$.}  Set
$u=\mu-2^L-3^L<3^L$ and pick $t,t'\in T_L$ with $u+t'=2t$
(Corollary~\ref{cor:reflect}).  Take
$\tau=3^L+t\in T_{L+1}$ (a leading ternary digit $1$;
Lemma~\ref{lem:digits}(i)) and $p=C_L+t'\in B_L$.  Then
\[
\mu+p \;=\; (u+t')+2\cdot 3^L+2^{L+1}+1
\;=\; 2(3^L+t)+2^{L+1}+1 \;=\; 2\tau+2^{L+1}+1 .
\]

\emph{Case D: $2^L+2\cdot 3^L\leq\mu$.}  Set
$\mu'=\mu-2\cdot 3^L-2^L<3^L+2^L$.  The induction hypothesis
gives $t\in T_L$, $p\in P_L$ with $\mu'+p=2t+2^L+1$; take
$\tau=3^L+t\in T_{L+1}$ and the same $p$:
\[
\mu+p \;=\; (2t+2^L+1)+2\cdot 3^L+2^L \;=\; 2\tau+2^{L+1}+1 .
\qedhere
\]
\end{proof}

\begin{proof}[Proof of Theorem~\ref{thm:blocked}]
Since $C_0=3\leq m$ and $C_L\to\infty$, there is a unique $L$
with $C_L\leq m<C_{L+1}$.  The interval $[C_L,C_{L+1})$ splits
into the \emph{span} $[C_L,\,C_L+3^L)$ of block $L$ and the
\emph{gap} $[C_L+3^L,\,C_{L+1})$, of length
$C_{L+1}-C_L-3^L=2^L+3^L$.

\emph{$m$ in the span.}  Write $m=C_L+u$ with $0\leq u<3^L$.
Since $m\notin\V$, $u\notin\Sd$.  Take
$x=C_L+\lo{u}$ and $y=C_L+\hi{u}$.  By Lemma~\ref{lem:witness},
$\lo{u},\hi{u}\in T_L$, so $x,y\in B_L\subseteq\V$; moreover
$\hi{u}<u$, so $y<m$ and $x\leq y<m$; and
$x+m=2C_L+u+\lo{u}=2C_L+2\hi{u}=2y$.  Finally $x\neq y$, since
$x=y$ would give $m=2y-x=y<m$.

\emph{$m$ in the gap.}  Write $m=C_L+3^L+\mu$ with
$0\leq\mu<2^L+3^L$, and take $t\in T_L$, $p\in P_L$ as in
Lemma~\ref{lem:gap}.  Set $x=p$ and $y=C_L+t\in B_L$.  Then
$x<C_L\leq y$, and $y\leq\max B_L<C_L+3^L\leq m$, and
\[
x+m \;=\; p+C_L+3^L+\mu \;=\; C_L+3^L+(2t+2^L+1)
\;=\; 2C_L+2t \;=\; 2y,
\]
using $C_L=2^L+3^L+1$.  (Note that in the gap the hypothesis
$m\notin\V$ was not needed: no term lies there, and this argument
re-derives that fact, since by Theorem~\ref{thm:admissible} a
blocked integer $>3$ cannot be a term.)
\end{proof}

\section{The greedy characterization}\label{sec:greedy}

We now assemble Theorems~\ref{thm:admissible}
and~\ref{thm:blocked} into the greedy recursion.  We take as the
formal definition of the sequence: $g(1)=1$, and for $n\geq 1$,
$g(n+1)$ is the least $k$ exceeding $g(1),\ldots,g(n)$ such that
$\{g(1),\ldots,g(n),k\}$ contains no $3$-term arithmetic
progression other than $(1,2,3)$.
The minimum exists: $k=2\max_i g(i)+1$ is always admissible,
since any new progression $x<y<k$ would need $2y=x+k>k$, while
every earlier term $y$ has $2y\leq 2\max_i g(i)<k$.  Note that
we do \emph{not} stipulate $g(2)=2$ and $g(3)=3$; the OEIS
definition does, but as the proof below shows (and as
Remark~\ref{rem:seed} isolates), the recursion started from the
single term $1$ derives the stipulated seed on its own.

\begin{theorem}\label{thm:greedy}
For every $n\geq 1$, $g(n)$ is the $n$-th smallest element of
$\V$.  In particular the term set of $g$ is $\V$
(Theorem~\ref{thm:blocks}), and Theorem~\ref{thm:main} follows.
\end{theorem}

\begin{proof}
Let $v_1<v_2<v_3<\cdots$ enumerate $\V$, so
$v_1,v_2,v_3,v_4,\ldots=1,2,3,6,\ldots$.  We prove
$g(n)=v_n$ by induction on $n$; the base $g(1)=1=v_1$ holds by
definition.

Assume $\{g(1),\ldots,g(n)\}=\{v_1,\ldots,v_n\}$.  The candidate
$k=v_{n+1}$ is admissible: it exceeds $v_n$, and
$\{v_1,\ldots,v_{n+1}\}\subseteq\V$, so by
Theorem~\ref{thm:admissible} the only progression it can contain
is the exempted $(1,2,3)$.

Conversely, let $v_n<k<v_{n+1}$; we show $k$ is inadmissible.
Since $v_1,v_2,v_3=1,2,3$ are consecutive integers, a $k$
strictly between $v_n$ and $v_{n+1}$ can exist only when
$v_n\geq 3$, so $k>3$.  Also $k\notin\V$, since the enumeration
has no elements between $v_n$ and $v_{n+1}$.  By
Theorem~\ref{thm:blocked} there are $x<y<k$ in $\V$ with
$x+k=2y$.  Since $x,y<k<v_{n+1}$, both lie in
$\{v_1,\ldots,v_n\}=\{g(1),\ldots,g(n)\}$.  Thus
$\{g(1),\ldots,g(n),k\}$ contains the progression $(x,y,k)$,
which is not $(1,2,3)$ because $k>3$.  Hence $k$ is
inadmissible, and $g(n+1)=v_{n+1}$.
\end{proof}

\begin{remark}[the seed is derived, not assumed]\label{rem:seed}
Unwinding the first steps: from $\{1\}$ the least admissible
extension is $2$ (no three terms exist), and from $\{1,2\}$ it is
$3$, because the unique progression in $\{1,2,3\}$ is the
exempted triple.  So the recursion reproduces
$a(1)=1$, $a(2)=2$, $a(3)=3$, and our formalization agrees with
the OEIS definition, which stipulates these values.  The
exemption is load-bearing: with $(1,2,3)$ forbidden instead of
exempted, the greedy sequence from $1$ is A003278
$=1,2,4,5,10,\ldots$, which differs already at the third term.
\end{remark}

\section{The summation form}\label{sec:sum}

The OEIS formula also expresses $s(n)$ as a sum over the $2$-adic
valuation $\nu_2$ (A007814); we record this with a proof, stated
division-free.  (The printed OEIS summand has a typographical
slip; see Appendix~\ref{app:errata}.)

\begin{proposition}\label{prop:sum}
For every $n\geq 0$,
\[
2\,s(n) \;=\; \sum_{k=1}^{n}\bigl(3^{\nu_2(k)}+1\bigr).
\]
Equivalently, $s(n)-s(n-1)=(3^{\nu_2(n)}+1)/2$ for $n\geq 1$.
\end{proposition}

\begin{proof}
It suffices to prove the increment form
$2\,s(n+1)=2\,s(n)+3^{\nu_2(n+1)}+1$ for $n\geq 0$ and sum; we
induct on $n$ (strongly).  If $n=2m$ is even, then
$s(n+1)=s(2m+1)=3s(m)+1=s(2m)+1=s(n)+1$ and $\nu_2(n+1)=0$, so
both sides increase by $3^0+1=2$.  If $n=2m+1$ is odd, then
$n+1=2(m+1)$, so $\nu_2(n+1)=\nu_2(m+1)+1$, and
\[
2\,s(n+1)-2\,s(n) \;=\; 6\,s(m+1)-6\,s(m)-2
\;=\; 3\bigl(3^{\nu_2(m+1)}+1\bigr)-2
\;=\; 3^{\nu_2(n+1)}+1,
\]
using the induction hypothesis at $m<n$.
\end{proof}

Combining with Theorem~\ref{thm:main}: for $n\geq 1$,
\[
a(n+2) \;=\; 1+2^{\lfloor\log_2 n\rfloor}
+\sum_{k=1}^{n}\frac{3^{\nu_2(k)}+1}{2},
\]
which is the OEIS formula with the summand read at index $k$.

\section{Not a Stanley sequence}\label{sec:stanley}

Given a finite $3$-AP-free set $A$, the \emph{Stanley sequence}
$S(A)$ of Odlyzko--Stanley \cite{OS1978} and
Erd\H{o}s--Lev--Rauzy--S\'andor--S\'ark\"ozy \cite{ELRSS1999} is
obtained by greedily appending to $A$ the least integer exceeding
$\max A$ that keeps the set $3$-AP-free.  The definition requires the
seed to be progression-free.  A092482 is \emph{not} of this form:
its seed $\{1,2,3\}$ is itself a $3$-term progression, and the
defining rule exempts, rather than avoids, that progression.  The
resemblance to Stanley sequences is nonetheless real --- by
Theorem~\ref{thm:blocks} the term set is a union of translated
initial segments of A003278$-1$ --- but it is blockwise, with the
translation constant $C_L=2^L+3^L+1$ changing at every power of
two.

This places A092482 outside the classification framework for
Stanley sequences.  Moy--Rolnick \cite{MoyRolnick2016} and
Rolnick \cite{Rolnick2017} organize Stanley sequences around the
\emph{regular} (Type-1) case, in which each successive
self-similar block of the sequence deviates from an exact
base-$3$ rescaling of the preceding ones by a \emph{constant},
the sequence's character.  For A092482 the deviation is not
constant: by Theorem~\ref{thm:blocks} the $L$-th block is
translated by $C_L=3^L+(2^L+1)$, so the correction to the pure
base-$3$ skeleton is $2^L+1$, which doubles (asymptotically) at
every power of two and is unbounded.  Equivalently, in the closed
form of Theorem~\ref{thm:main} the correction term
$2^{\lfloor\log_2 n\rfloor}$ is unbounded, where a regular
Stanley sequence would exhibit a constant character.  So the
classification machinery does not apply, and A092482 does not
appear among its examples.

A nearby OEIS structure also fails to contain it.  The array
A093682 collects greedy nonarithmetic-$3$-progression sequences
with conjectured closed forms; its comments state, verbatim:

\begin{quote}
\texttt{The nonarithmetic-3-progression sequences starting with
a(1)=1, a(2)=1+3\textasciicircum m or 1+2*3\textasciicircum m, m
>= 0, seem to have especially simple 'closed' forms. None of
these formulas have been proved, however.}
\end{quote}

\noindent Its rows are seeded $1,\,1+3^m$ or $1,\,1+2\cdot 3^m$
(rows $0$--$6$ are A003278, A004793, A033157, and
A093678--A093681), and the conjectured closed forms there carry a
\emph{$P$-periodic} correction term ($T(m,n)=\sum_{k=1}^{n-1}
(3^{A007814(k)}+1)/2+f(n)$ with $f$ a $P$-periodic function).
A092482 is not a row of A093682: its seed does not have the
required shape, and its correction $2^{\lfloor\log_2 n\rfloor}$
is unbounded rather than periodic.  (A093682 cross-references
A092482; the theorem proved here settles a conjecture in that
circle which lies outside the A093682 family.  The conjectured
formulas for the rows of A093682 itself remain open.)

\section{Machine verification}\label{sec:lean}

The proof in this note has been formalized in Lean~4
\cite{Lean4} over the Mathlib library \cite{mathlib}.  The
formalization is available at \url{TODO}.  Its headline statement,
\texttt{greedySeq\_eq\_closedForm}, asserts that the greedy
sequence equals the closed form as functions on $\N$; the
development compiles with no \texttt{sorry} and no errors, and
Lean's axiom audit reports exactly the three standard axioms
\texttt{propext}, \texttt{Classical.choice}, and
\texttt{Quot.sound}.

Three design points bear on what the verification guarantees.
First, there is no circularity: the greedy sequence is defined in
Lean directly from the OEIS rule --- the next term is obtained by
\texttt{Nat.find}, the least $k$ exceeding all previous terms
whose insertion creates no progression other than $(1,2,3)$ ---
with no reference to the closed form; the closed form is defined
separately, and the theorem connects the two.  Second, the seed
is derived, not assumed: the formal recursion starts from the
singleton $\{1\}$ and proves that its second and third values are
$2$ and $3$ (Remark~\ref{rem:seed}), so the formal definition
matches the OEIS definition, which stipulates those values.
Third, the first $57$ terms of the formal greedy sequence are
proved equal, by a kernel-checked computation, to the $57$ terms
printed in the OEIS entry's data field; no use is made of
\texttt{native\_decide} (i.e., no compiled-code oracle), so the
entire development, including the ground computation, is checked
by the Lean kernel.  The blockwise structure of
Sections~\ref{sec:blocks}--\ref{sec:greedy} mirrors the formal
proof, which reuses a formalization of the classical A003278
digit argument.

\appendix

\section{Errata for the OEIS entry}\label{app:errata}

The formula field of A092482 quoted in
Section~\ref{sec:intro} contains two typographical slips.  Both
are confirmed by direct computation against the entry's own data
field, and both are repaired by the readings proved in this note.

\begin{enumerate}
\item \emph{The summand index.}  The sum is printed as
\texttt{Sum\_\{k=1..n\} (3\textasciicircum A007814(n) + 1)/2}:
the summand must depend on the summation index $k$, not on $n$.  Read
literally, the sum equals $n\,(3^{\nu_2(n)}+1)/2$, which at $n=3$
gives $3\cdot(3^0+1)/2=3$ and hence $a(5)=1+2+3=6$, whereas the
data field has $a(5)=7$.  With summand
$(3^{\nu_2(k)}+1)/2$ the sum equals $s(n)$
(Proposition~\ref{prop:sum}), as required.
\item \emph{Indexing of the cross-references.}  In
$a(n+2)=1+A053644(n)+A005836(n)$, both cross-references must be
read as $0$-indexed lookups into the printed term lists, i.e.,
$A053644(n)=2^{\lfloor\log_2 n\rfloor}$ (from the list
$0,1,2,2,4,4,4,4,8,\ldots$) and $A005836(n)=s(n)$ (from the list
$0,1,3,4,9,10,12,13,27,\ldots$).  Reading A005836 as $1$-indexed
(its declared OEIS offset) gives $A005836(3)=3$ and hence
$a(5)=1+2+3=6\neq 7$; reading A053644 as $1$-indexed gives
$A053644(4)=2$ and hence $a(6)=1+2+9=12\neq 14$.
\end{enumerate}

\noindent Neither slip affects the intended formula, which the
entry's Mathematica program (Alcover, 2019) renders unambiguously
and which Theorem~\ref{thm:main} proves.

\begin{thebibliography}{99}

\bibitem{OEIS}
OEIS Foundation Inc., \emph{The On-Line Encyclopedia of Integer
Sequences}, 2026.  Available at \url{https://oeis.org}.  Entries
A092482, A003278, A005836, A007814, A053644, A093682.

\bibitem{OS1978}
A.~M. Odlyzko and R.~P. Stanley, \emph{Some curious sequences
constructed with the greedy algorithm}, Bell Laboratories
internal memorandum, 1978.

\bibitem{ELRSS1999}
P.~Erd\H{o}s, V.~Lev, G.~Rauzy, C.~S\'andor, and A.~S\'ark\"ozy,
Greedy algorithm, arithmetic progressions, subset sums and
divisibility, \emph{Discrete Math.} \textbf{200} (1999),
119--135.

\bibitem{MoyRolnick2016}
R.~A. Moy and D.~Rolnick, Novel structures in Stanley sequences,
\emph{Discrete Math.} \textbf{339} (2016), 689--698.

\bibitem{Rolnick2017}
D.~Rolnick, On the classification of Stanley sequences,
\emph{European J. Combin.} \textbf{59} (2017), 51--70.

\bibitem{Fried2024}
S.~Fried, Proofs of some conjectures from the OEIS, preprint,
2024.  Available at \url{https://arxiv.org/abs/2410.07237}.

\bibitem{Fried2026}
S.~Fried, Further proofs of conjectures from the OEIS, preprint,
2026.  Available at \url{https://arxiv.org/abs/2607.24832}.

\bibitem{Lean4}
L.~de~Moura and S.~Ullrich, The Lean 4 theorem prover and
programming language, in \emph{Automated Deduction --- CADE 28},
Lecture Notes in Comput.\ Sci., Vol.~12699, Springer, 2021,
pp.~625--635.

\bibitem{mathlib}
The mathlib Community, The Lean mathematical library, in
\emph{Proceedings of the 9th ACM SIGPLAN International Conference
on Certified Programs and Proofs (CPP 2020)}, ACM, 2020,
pp.~367--381.

\end{thebibliography}

\end{document}
